%%%%%%%%%%  scan idx 20  |  folio 7  %%%%%%%%%%
auxquels on ne saurait associer de classe dans $K^\bullet(Y)$, cependant le complexe $Rf_*(\underline F)$ sur $Y$, image directe de $\underline F$ au sens des catégories dérivées, dont les $R^if_*(\underline F)$ sont des faisceaux de cohomologie, a tendance à être excellent. De façon précise, supposons que $\underline F$ soit de tor-dimension finie sur $\underline O_Y$, ce qui sera le cas si $\underline O_X$ est de tor-dimension finie sur $\underline O_Y$ et si de plus $F$ est localement libre. Alors on montre (Exp.\ III) que le complexe $Rf_*(F)$ est «~\underline{parfait}~», par quoi on entend ici~: à faisceaux de cohomologie cohérents (ce qu'on savait déjà de toutes façons) \underline{et} de tor-dimension globale finie~; ou encore~: sur tout ouvert affine $U$ de $Y$, ce complexe est isomorphe, dans la catégorie dérivée $D(U)$, à un complexe à degrés bornés et localement libre de type fini en chaque degré. Du point de vue de l'Algèbre Homologique, les complexes parfaits sont «~aussi bons~» que les Modules localement libres, dont ils constituent la généralisation naturelle. Ils forment d'ailleurs une «~sous-catégorie triangulée~» de la catégorie dérivée $D(Y)$, i.e.\ une sous-catégorie stable par les cônes (ou «~mapping cylinders~») de morphismes. D'autre part, à toute catégorie triangulée $C$ on associe un groupe $K(C)$ de classes, variante naturelle de la construction du groupe des classes d'une catégorie abélienne, en prenant le groupe libre engendré par les objets de $C$, divisé par les relations
\[
\cl(K)-\cl(K')-\cl(K'')
\]
pour tout triangle distingué
\[
\begin{tikzcd}[column sep=large]
 & K'' \arrow[dl] & \\
K' \arrow[rr] & & K \arrow[ul] .
\end{tikzcd}
\]
Ce qui remplace alors l'observation fondamentale de [RRR] dans le cas actuel est le fait que, lorsque le schéma $Y$ admet un Module inversible ample, tout complexe parfait $L_\bullet$ sur $Y$ est globalement (et non seulement localement) isomorphe dans $D(Y)$ à un complexe $L'_\bullet$ à degrés bornés et localement libre de type fini en chaque degré (Exp.\ II)~; d'où il résulte facilement (Exp.\ IV) que l'homomorphisme canonique
\[
K^\bullet(Y)\longrightarrow K(\Parf(Y)),
\]
où $\Parf(Y)$ est la catégorie triangulée des complexes parfaits sur $Y$, est un isomorphisme. Donc à tout complexe parfait $L_\bullet$ sur $Y$ on peut associer une classe
\[
\cl^\bullet(L_\bullet)\in K^\bullet(Y),
\]
